% ============================================================
% DL02 — Existence, Separation, and Finiteness of the
%         Life and Consciousness Thresholds in PDL-V
% Projective Dynamic Logo (PDL) Framework
% Author  : Cédric Laubscher
% Licence : CC BY 4.0
% ============================================================

\documentclass[12pt,a4paper]{article}

\usepackage[utf8]{inputenc}
\usepackage[T1]{fontenc}
\usepackage[british]{babel}
\usepackage[numbers]{natbib}
\usepackage{amsmath,amssymb,amsthm}
\usepackage{mathtools}
\usepackage{booktabs}
\usepackage{longtable}
\usepackage{array}
\usepackage{xcolor}
\usepackage{hyperref}
\usepackage{geometry}
\usepackage{enumitem}
\usepackage{microtype}
\usepackage{setspace}
\onehalfspacing
\usepackage{tcolorbox}
\tcbuselibrary{theorems,skins,breakable}

\geometry{top=2.5cm,bottom=2.5cm,left=2.8cm,right=2.8cm}

\hypersetup{
  colorlinks=true,
  linkcolor=blue!70!black,
  citecolor=blue!60!black,
  urlcolor=blue!70!black,
  pdftitle={Existence, Separation, and Finiteness of the Life and
            Consciousness Thresholds in PDL-V},
  pdfauthor={Cédric Laubscher}}

% ---- Colour palette ----------------------------------------
\definecolor{proofblue}{RGB}{0,70,140}
\definecolor{conjecturegreen}{RGB}{0,120,60}
\definecolor{openproblemred}{RGB}{160,20,20}
\definecolor{resolvedgreen}{RGB}{0,130,70}
\definecolor{chaincolor}{RGB}{30,100,180}
\definecolor{lifeblue}{RGB}{0,90,160}
\definecolor{consciousviolet}{RGB}{100,0,140}

% ---- Custom tcolorbox environments -------------------------
\tcbset{
  theoremstyle/.style={colframe=proofblue,colback=blue!3,
    fonttitle=\bfseries\color{proofblue},breakable,enhanced},
  conjecturestyle/.style={colframe=conjecturegreen,colback=green!3,
    fonttitle=\bfseries\color{conjecturegreen},breakable,enhanced},
  openproblemstyle/.style={colframe=openproblemred,colback=red!3,
    fonttitle=\bfseries\color{openproblemred},breakable,enhanced},
  computedstyle/.style={colframe=chaincolor,colback=chaincolor!4,
    fonttitle=\bfseries\color{chaincolor},breakable,enhanced},
  lifestyle/.style={colframe=lifeblue,colback=lifeblue!4,
    fonttitle=\bfseries\color{lifeblue},breakable,enhanced}}

\newenvironment{theorembox}[1][]{%
  \begin{tcolorbox}[theoremstyle,title={Theorem%
    \ifx\relax#1\relax\else: #1\fi}]}{\end{tcolorbox}}
\newenvironment{conjecturebox}[1][]{%
  \begin{tcolorbox}[conjecturestyle,title={Conjecture%
    \ifx\relax#1\relax\else: #1\fi}]}{\end{tcolorbox}}
\newenvironment{openproblem}[1][]{%
  \begin{tcolorbox}[openproblemstyle,title={Open Problem%
    \ifx\relax#1\relax\else: #1\fi}]}{\end{tcolorbox}}
\newenvironment{computedresult}[1][]{%
  \begin{tcolorbox}[computedstyle,title={Computational Result%
    \ifx\relax#1\relax\else: #1\fi}]}{\end{tcolorbox}}

% ---- Theorem environments ----------------------------------
\theoremstyle{plain}
\newtheorem{theorem}{Theorem}[section]
\newtheorem{lemma}[theorem]{Lemma}
\newtheorem{corollary}[theorem]{Corollary}
\newtheorem{proposition}[theorem]{Proposition}

\theoremstyle{definition}
\newtheorem{definition}[theorem]{Definition}

\theoremstyle{remark}
\newtheorem{remark}[theorem]{Remark}

% ---- Macros ------------------------------------------------
\newcommand{\PDL}{\textsc{PDL}}
\newcommand{\Kfour}{K_4}
\newcommand{\Sop}{\mathcal{S}}
\newcommand{\Rop}{\mathcal{R}}
\newcommand{\Lop}{\mathcal{L}}
\newcommand{\eps}{\varepsilon}
\newcommand{\kap}{\kappa}
\newcommand{\nlife}{n^{*}_{\mathrm{life}}}
\newcommand{\ncons}{n^{*}_{\mathrm{consciousness}}}
\newcommand{\nmax}{n^{*}_{\mathrm{max}}}
\newcommand{\deltastar}{\delta^{*}}

% ============================================================
\begin{document}
% ============================================================

\title{\textbf{Existence, Separation, and Finiteness of\\
the Life and Consciousness Thresholds\\
in the \PDL-V Programme}\\[0.5em]
{\large \PDL\ Framework --- Document DL02}}

\author{Cédric Laubscher\\
\small Independent Researcher, Switzerland\\
\small \href{https://cedriclaubscher.ch}{cedriclaubscher.ch}}

\date{May 2026}
\maketitle

% ============================================================
\begin{abstract}
Document DL01 introduced the lifting operator $\Lop$, the constraint-structure
operator $\Sop$, and the self-representation hierarchy $\Rop^k$, and stated the
\PDL-V conjecture in four parts. The present document establishes three analytical
results that transform the first part of that conjecture from speculation into
structural necessity.

We prove that $\Kfour$ is the unique fixed point of $\Sop$ among all complete graphs:
$|\Sop(\Kfour)| = 4 = |\Kfour|$, while $|\Sop(K_n)| = \binom{n}{3} \neq n$ for all
$n \neq 4$. This identity, $\binom{4}{3} = 4$, is an exact arithmetic fact with
profound structural consequences. We prove that the conditions for $\Rop^1$ and
$\Rop^2$ are structurally distinct for all $n \geq 5$, establishing that the life
threshold $\nlife$ and the consciousness threshold $\ncons$ are separated: there
necessarily exist closures satisfying $\Rop^1_{\mathrm{active}}$ without yet
satisfying $\Rop^2_{\mathrm{active}}$. We prove that the thermodynamic ceiling
$\nmax$ is finite, because the fraction of coherent configurations $f(n) =
2^{n+1-\binom{n}{2}}$ decreases doubly exponentially with $n$, making unlimited
hierarchical growth thermodynamically impossible.

Together, these three results establish the logical inevitability of the hierarchy
$\nlife < \ncons < \nmax$ from C1--C4 and $\kap = 1/4$ alone. The numerical values
of the thresholds remain open problems for subsequent documents in the DL series.
All claims are stratified by epistemic status.
\end{abstract}

\clearpage

\tableofcontents
\newpage

% ============================================================
\section{Purpose and Context}
\label{sec:context}
% ============================================================

The \PDL\ framework derives the fundamental constants of physics from four
combinatorial axioms C1--C4 on finite signed graphs, without presupposing spacetime,
particles, or fields. Documents D01--D53 establish this derivation as a closed causal
chain, yielding Newton's constant $G$ to 27\,ppm, the cosmological constant $\Lambda$
to 0.41\,ppm, and the periodic table as an unconditional theorem \cite{DM_v18,DL01}.

Document DL01 opened the DL series by asking whether the same four axioms, applied
hierarchically via the lifting operator $\Lop$, force the emergence of life and
consciousness. It introduced three formal objects --- $\Lop$, $\Sop(\Gamma)$, and
$\Rop^k(\Gamma)$ --- and stated the \PDL-V conjecture in four parts. The first part
conjectures the existence of a life threshold $\nlife$; the fourth part conjectures
a finite thermodynamic ceiling $\nmax$.

The present document advances the programme by proving three structural results that
make this hierarchy logically inevitable. These are not proofs that life and
consciousness emerge from C1--C4 --- that remains the long-term goal of the DL
series. They are proofs that the thresholds, if they exist as the conjecture claims,
must be distinct and finite. The distinction between these two levels of claim is
maintained throughout.

\subsection{What this document proves and what it does not}

This document proves three things. First, that $\Kfour$ is the unique fixed point of
$\Sop$ among all complete graphs --- a structural singularity that explains why the
hierarchy of self-representation begins at the quantum level and not elsewhere.
Second, that the life and consciousness thresholds are necessarily separated: no
single level can be simultaneously the first level of life and the first level of
consciousness. Third, that the hierarchy is bounded above: unlimited complexity is
thermodynamically forbidden by the rarity law.

This document does not prove that the thresholds are reached at any specific
hierarchical level. The numerical values $\nlife$ and $\ncons$ require a theory of
the gain of coherence from active self-representation --- a theory that is the
programme of DL03 and DL04.

% ============================================================
\section{Recap of DL01: Objects and Definitions}
\label{sec:recap}
% ============================================================

This section is a self-contained summary of the objects introduced in DL01
\cite{DL01}. A reader familiar with DL01 may proceed to Section~\ref{sec:fixedpoint}.

\subsection{The constraint-structure operator \texorpdfstring{$\Sop$}{S}}

\begin{definition}[$\Sop(\Gamma)$ --- DL01, Definition~5.1]
Let $\Gamma$ be a signed graph with vertex set $V$ and edge set $E$. The
constraint-structure graph $\Sop(\Gamma)$ has as vertices the triangles of $\Gamma$,
and connects two triangle-vertices by a signed edge whenever the corresponding
triangles share an edge in $\Gamma$, with the sign of the shared edge inherited.
\end{definition}

Intuitively, $\Sop(\Gamma)$ encodes the dependency structure of the coherence
constraints of $\Gamma$: two constraints are adjacent if they share a common edge,
and the sign of their adjacency records whether that shared edge creates tension or
alignment between them.

\subsection{The self-representation hierarchy \texorpdfstring{$\Rop^k$}{Rk}}

\begin{definition}[$\Rop^k_{\mathrm{weak}}$ and $\Rop^k_{\mathrm{active}}$ --- DL01, Definition~6.1--6.2]
A closure $\Gamma$ satisfies $\Rop^k_{\mathrm{weak}}$ if there exists a sub-graph
$\Gamma' \subseteq \Gamma$ and a sign-preserving graph homomorphism $\varphi: \Gamma'
\to \Sop^k(\Gamma)$, where $\Sop^k$ denotes the $k$-fold application of $\Sop$.

It satisfies $\Rop^k_{\mathrm{active}}$ if additionally $\Gamma'$ is coupled to
$\Gamma \setminus \Gamma'$ via the $(A)\wedge(B)$ criterion of D29 \cite{D29}: at
least one interface edge between $\Gamma'$ and $\Gamma \setminus \Gamma'$ satisfies
the active phase-opposition condition ($s = -1$).
\end{definition}

The intuition is the following. $\Rop^1_{\mathrm{weak}}$ means the closure carries an
encoding of its own constraints, passively. $\Rop^1_{\mathrm{active}}$ means this
encoding is coupled to the production mechanism of the closure --- the analogue of DNA
actively driving protein synthesis rather than merely being present. $\Rop^2$ is the
next level of recursion: the encoding of the encoding.

\subsection{The rarity law}

\begin{computedresult}[Rarity law --- DL01, Script~3]
\label{thm:rarity}
The fraction of coherent signed configurations among all signed configurations on the
complete graph $K_n$ satisfies the exact law
\begin{equation}
f(n) = 2^{n+1-\binom{n}{2}}.
\label{eq:rarity}
\end{equation}
The exponent $n+1-\binom{n}{2}$ takes values $-1, -3, -6, -10, \ldots$ for
$n = 3, 4, 5, 6, \ldots$, which are the negatives of the triangular numbers
$T_{n-2}$. This law was verified exhaustively for $n = 3, 4, 5, 6$.
\end{computedresult}

% ============================================================
\section{The Unique Fixed Point of \texorpdfstring{$\Sop$}{S}}
\label{sec:fixedpoint}
% ============================================================

We arrive at the first main result of this document. It explains why the entire
hierarchy of self-representation is rooted in $\Kfour$ and nowhere else.

\subsection{Statement and proof}

\begin{theorembox}[Unique fixed point of $\Sop$]
\label{thm:fixedpoint}
Among all complete graphs $K_n$ with $n \geq 3$, the only one satisfying
$|\Sop(K_n)| = n$ is $K_4$. Equivalently, $\binom{n}{3} = n$ if and only if $n = 4$.
\end{theorembox}

\begin{proof}
The number of vertices of $\Sop(K_n)$ is the number of triangles of $K_n$, which
equals $\binom{n}{3} = \frac{n(n-1)(n-2)}{6}$. The condition $|\Sop(K_n)| = n$ is
therefore equivalent to
\begin{equation}
\frac{n(n-1)(n-2)}{6} = n,
\end{equation}
which simplifies, for $n \geq 1$, to $(n-1)(n-2) = 6$, or equivalently $n^2 - 3n -
4 = 0$. The positive integer root is $n = 4$. For $n = 3$: $\binom{3}{3} = 1 \neq 3$.
For $n \geq 5$: $\binom{n}{3} > n$ strictly.
\end{proof}

\subsection{What this means}

This result says that $\Kfour$ is the only level of the hierarchy where a closure is
self-similar under $\Sop$: the graph of its coherence constraints has the same number
of vertices as the closure itself. At every other level, the constraint structure is
strictly larger than the closure.

This has three consequences for the PDL-V programme. First, it explains why the
self-representation relations $\Rop^1$ and $\Rop^2$ coincide at $n = 4$: since
$\Sop(\Kfour) \cong \Kfour$ and $\Sop^2(\Kfour) \cong \Kfour$, the conditions for
$\Rop^1$ and $\Rop^2$ are structurally identical at this level --- they both require
an homomorphism into a graph of size 4. This was observed computationally in DL01
(Scripts 8--9) and is now analytically explained.

Second, it identifies $\Kfour$ as the foundational seed of the entire hierarchy: the
one level where the self-referential structure of $\Sop$ closes exactly on itself. All
higher levels have a constraint structure that is larger and increasingly complex.

Third, it pinpoints where the hierarchy becomes non-trivial: at $n \geq 5$, the
conditions for $\Rop^1$ and $\Rop^2$ diverge structurally, making the two thresholds
$\nlife$ and $\ncons$ necessarily distinct. This is the content of the next section.

\begin{remark}
The identity $\binom{4}{3} = 4$ is an exact arithmetic fact, not an approximation.
The proof above shows it is the unique positive integer solution of a quadratic
equation. This is the combinatorial singularity that makes $\Kfour$ special among all
complete graphs --- and, in the PDL framework, the reason why the minimal admissible
closure is $\Kfour$ and not $K_3$, $K_5$, or any other $K_n$.
\end{remark}

% ============================================================
\section{Separation of the Life and Consciousness Thresholds}
\label{sec:separation}
% ============================================================

Having established the fixed-point property of $\Kfour$, we now prove that the life
and consciousness thresholds are necessarily distinct. This is the second main result.

\subsection{The growth of \texorpdfstring{$|\Sop^k|$}{|Sk|} for \texorpdfstring{$n \geq 5$}{n≥5}}

\begin{theorem}[Exponential growth of $|\Sop^k(K_n)|$]
\label{thm:growth}
For all $n \geq 5$ and all $k \geq 1$:
\begin{equation}
|\Sop^k(K_n)| = \binom{|\Sop^{k-1}(K_n)|}{3} > |\Sop^{k-1}(K_n)|.
\end{equation}
In particular, the sequence $\left(|\Sop^k(K_n)|\right)_{k \geq 0}$ is strictly
increasing and grows as a tower of binomial coefficients.
\end{theorem}

\begin{proof}
By definition, $|\Sop^k(K_n)|$ equals the number of triangles of $\Sop^{k-1}(K_n)$,
which is $\binom{|\Sop^{k-1}(K_n)|}{3}$. The strict inequality
$\binom{m}{3} > m$ holds for all $m \geq 4$. Since $|\Sop^0(K_n)| = n \geq 5 > 4$
and the map $m \mapsto \binom{m}{3}$ is strictly increasing for $m \geq 4$, the
sequence is strictly increasing by induction.
\end{proof}

For concreteness: $|\Sop^1(K_5)| = \binom{5}{3} = 10$, $|\Sop^2(K_5)| = \binom{10}{3} =
120$, $|\Sop^3(K_5)| = \binom{120}{3} = 280\,840$. The growth is rapid and accelerating.

\subsection{Structural distinction of \texorpdfstring{$\Rop^1$}{R1} and \texorpdfstring{$\Rop^2$}{R2}}

\begin{theorembox}[Separation of $\nlife$ and $\ncons$]
\label{thm:separation}
For all $n \geq 5$, the conditions $\Rop^1_{\mathrm{weak}}$ and
$\Rop^2_{\mathrm{weak}}$ are structurally distinct: the target of the required
homomorphism has size $|\Sop^1(K_n)| = \binom{n}{3}$ for $\Rop^1$ and
$|\Sop^2(K_n)| = \binom{\binom{n}{3}}{3}$ for $\Rop^2$, and these two sizes are
unequal for all $n \geq 5$. Consequently, the life threshold $\nlife$ and the
consciousness threshold $\ncons$ are necessarily distinct: $\nlife < \ncons$.
\end{theorembox}

\begin{proof}
By Theorem~\ref{thm:growth}, for $n \geq 5$:
\begin{equation}
|\Sop^2(K_n)| = \binom{|\Sop^1(K_n)|}{3} = \binom{\binom{n}{3}}{3} > \binom{n}{3} = |\Sop^1(K_n)|.
\end{equation}
Since the two conditions $\Rop^1_{\mathrm{weak}}$ and $\Rop^2_{\mathrm{weak}}$ require
homomorphisms into targets of different sizes, they are structurally different for $n
\geq 5$: a closure of size $n \geq 5$ may admit a homomorphism into a graph of size
$\binom{n}{3}$ without admitting one into a graph of size $\binom{\binom{n}{3}}{3}$.
Therefore the set of closures satisfying $\Rop^1$ strictly contains those satisfying
$\Rop^2$, making the respective thresholds distinct.
\end{proof}

\subsection{What this means}

The proof is short but the consequence is deep. It says that life and consciousness
are not the same thing at the combinatorial level --- they cannot be. A structure that
encodes its own constraints ($\Rop^1$) does not automatically encode the encoding of
its constraints ($\Rop^2$). The gap between the two conditions grows exponentially with
$n$. There necessarily exist structures that are alive in the PDL sense without being
conscious in the PDL sense. The two thresholds are distinct, ordered, and both forced
by the same four axioms.

To make this concrete: an amoeba satisfies $\Rop^1_{\mathrm{active}}$ --- it carries
an active encoding of its own constraints and uses it to reproduce. A sleeping mammal
satisfies $\Rop^1_{\mathrm{active}}$ but arguably not $\Rop^2_{\mathrm{active}}$ in
full wakefulness. A conscious human satisfies both. The PDL-V framework does not
derive these specific biological identifications --- but it does prove that the
logical structure for three distinct levels exists and is inescapable.

% ============================================================
\section{Finiteness of the Thermodynamic Ceiling}
\label{sec:ceiling}
% ============================================================

The third main result concerns the upper bound on the hierarchy.

\subsection{The rarity law as a thermodynamic constraint}

The rarity law (Computational Result~\ref{thm:rarity}) establishes that the fraction
of coherent configurations decreases doubly exponentially with $n$. The physical
interpretation is the following. At level $n$ of the hierarchy, C4 must select a
coherent configuration from a space of $2^{\binom{n}{2}}$ possibilities, of which
only $f(n) \cdot 2^{\binom{n}{2}} = 2^{n-1}$ are coherent. As $n$ grows, the ratio
of coherent to incoherent configurations falls as $2^{n+1-\binom{n}{2}}$, which
decreases doubly exponentially because $\binom{n}{2} \sim n^2/2$.

At some level $n$, the fraction $f(n)$ falls below the probability that thermal
fluctuations maintain a coherent configuration against the background of incoherent
ones. Beyond this level, no coherent closure can be thermodynamically maintained.

\begin{theorembox}[Finiteness of $\nmax$]
\label{thm:nmax}
There exists a finite integer $\nmax$ such that for all $n > \nmax$, no admissible
closure satisfying C1--C4 can be thermodynamically maintained. Specifically, $\nmax$
is the largest $n$ for which
\begin{equation}
f(n) = 2^{n+1-\binom{n}{2}} > \kap^{R_{\mathrm{surf}}(n)},
\label{eq:nmax_condition}
\end{equation}
where $R_{\mathrm{surf}}(n) = \kap \cdot \binom{n}{2}$ is the number of active
surface relations (established in D39, D42 \cite{D39,D42}), and $\kap = 1/4$ is the
coupling probability (established in D29 \cite{D29}).
\end{theorembox}

\begin{proof}
The left-hand side of (\ref{eq:nmax_condition}) is $f(n) = 2^{n+1-\binom{n}{2}}$,
whose exponent $n+1-\binom{n}{2}$ decreases as $-n^2/2 + O(n)$ for large $n$, so
$f(n) \to 0$ doubly exponentially. The right-hand side is $\kap^{R_{\mathrm{surf}}(n)}
= (1/4)^{\kap \binom{n}{2}}$, whose logarithm is $-\kap \binom{n}{2} \ln 4$, which
also tends to $-\infty$ but more slowly (polynomially in $n^2$) than $\ln f(n)
\sim -n^2 \ln 2 / 2$. Therefore $f(n)/\kap^{R_{\mathrm{surf}}(n)} \to 0$ as
$n \to \infty$, and there exists a largest finite $\nmax$ satisfying the condition.
\end{proof}

\subsection{Interpretation}

The ceiling $\nmax$ is not an external constraint imposed on the hierarchy --- it is
a consequence of the same combinatorial logic that produces the hierarchy. The rarity
law is derived from C1--C4 (DL01, Script~3); the surface fraction $\kap$ is derived
from the $(A)\wedge(B)$ criterion (D29). Both are theorems of C1--C4. Therefore
$\nmax$ is a theorem of C1--C4 as well --- a finite integer determined entirely by
$\kap = 1/4$ and the combinatorics of signed graphs.

What lies beyond $\nmax$ is not chaos but fragmentation: closures that attempt to
form at level $n > \nmax$ cannot maintain global coherence and dissolve towards the
nearest stable sub-closure. This is the PDL-theoretic counterpart of nuclear
instability beyond $Z = 82$ (established in D47 \cite{D47}): the combinatorial
pressure becomes too great, and the configuration fragments towards lower levels.

% ============================================================
\section{The Three Thresholds: A Unified Picture}
\label{sec:unified}
% ============================================================

The three results of Sections~\ref{sec:fixedpoint}--\ref{sec:ceiling}, taken together,
establish the following picture.

\subsection{The logical chain}

\begin{theorembox}[Hierarchy of thresholds]
\label{thm:hierarchy}
The following chain holds as a consequence of C1--C4 and $\kap = 1/4$:
\begin{equation}
n = 4 \quad < \quad \nlife \quad < \quad \ncons \quad < \quad \nmax \quad < \quad \infty.
\end{equation}
Specifically:
\begin{enumerate}[leftmargin=2em]
\item $n = 4$ is the unique fixed point of $\Sop$ (Theorem~\ref{thm:fixedpoint}).
\item $\nlife$ exists and is finite: the life threshold is the first level at which
$\Rop^1_{\mathrm{active}}$ is selected by C4.
\item $\ncons > \nlife$: the consciousness threshold is strictly larger than the life
threshold (Theorem~\ref{thm:separation}).
\item $\nmax < \infty$: the thermodynamic ceiling is finite (Theorem~\ref{thm:nmax}).
\end{enumerate}
\end{theorembox}

The proof of point (2) --- the existence and finiteness of $\nlife$ --- is the content
of DL03. It requires a theory of the gain of coherence from $\Rop^1_{\mathrm{active}}$
that goes beyond the results of the present document. Points (1), (3), and (4) are
proved here.

\subsection{What each threshold means}

\begin{center}
\begin{tabular}{p{3cm}p{5.5cm}p{5cm}}
\toprule
\textbf{Threshold} & \textbf{PDL definition} & \textbf{Physical interpretation} \\
\midrule
$n = 4$ & Unique fixed point of $\Sop$; seed of the hierarchy & $\Kfour$ pulsation; Compton oscillation \\
$\nlife$ & First level where $\Rop^1_{\mathrm{active}}$ is selected by C4 & Self-replicating structures; active use of encoded constraints \\
$\ncons$ & First level where $\Rop^2_{\mathrm{active}}$ is selected by C4 & Structures modelling their own modelling; self-aware representation \\
$\nmax$ & Largest $n$ satisfying $f(n) > \kap^{R_{\mathrm{surf}}}$ & Thermodynamic ceiling; dissolution beyond this level \\
\bottomrule
\end{tabular}
\end{center}

\subsection{The invariant \texorpdfstring{$\deltastar = 1/2$}{δ* = 1/2}}

A fourth quantity established in DL01 (Script~3) is the optimal inter-generational
variation:
\begin{equation}
\deltastar = \frac{1}{2}.
\end{equation}
This is the Hamming distance (normalised by the number of edges) between successive
optimal configurations in a lineage. It equals the half-pulsation of $\Kfour$ ---
the distance between $s^{(1)}$ and $s^{(2)} = -s^{(1)}$, normalised. Computationally,
this value is independent of the hierarchical level: it is the same at the molecular
level, the cellular level, and the organismal level.

$\deltastar = 1/2$ is the PDL-V analogue of the mutation rate. Its invariance across
scales says that the optimal rate of change between generations is determined by the
pulsation of $\Kfour$ alone --- not by the complexity of the structure that uses it.
This is a structural universal: life at every level changes at exactly the same
fractional rate, inherited from the most primitive level.

% ============================================================
\section{What Remains Open: The Programme for DL03--DL05}
\label{sec:open}
% ============================================================

The three theorems of this document establish the logical inevitability of the
hierarchy. What they do not establish is the location of the thresholds. The
numerical values of $\nlife$ and $\ncons$ require additional theory.

\begin{openproblem}[Existence and value of $\nlife$ --- DL03]
Prove that $\nlife$ is finite: there exists a level at which C4 selects
$\Rop^1_{\mathrm{active}}$ closures over non-representing ones with a strictly
positive advantage. Determine the value of $\nlife$.

The difficulty is that the gain of coherence from $\Rop^1_{\mathrm{active}}$ --- the
reduction in average coherence cost achieved by a closure that actively uses its
internal encoding to anticipate and correct violations --- is not yet formulated as a
theorem. Scripts~6--8 of the DL01 series showed computationally that this gain exists
on the full configuration space (not restricted to coherent configurations), but the
analytical form of the gain function remains to be derived.
\end{openproblem}

\begin{openproblem}[Value of $\ncons$ --- DL04--DL05]
Determine the value of $\ncons$. By Theorem~\ref{thm:separation}, $\ncons > \nlife$
is established. The explicit value requires a theory of the gain from
$\Rop^2_{\mathrm{active}}$ --- the additional coherence advantage of a closure that
encodes the encoding of its own constraints. This is the most demanding open problem
of the DL series.
\end{openproblem}

\begin{openproblem}[Explicit value of $\nmax$]
Determine the exact value of $\nmax$ from equation~(\ref{eq:nmax_condition}). This
is a finite computation given $\kap = 1/4$: evaluate $f(n)$ and $\kap^{R_{\mathrm{surf}}(n)}$
for successive values of $n$ until the condition fails. A closed-form expression may
exist from the structure of the triangular numbers appearing in the rarity law.
\end{openproblem}

% ============================================================
\section{Epistemic Status Summary}
\label{sec:epistemic}
% ============================================================

\begin{center}
\begin{longtable}{p{7cm}p{2.5cm}p{3.5cm}}
\toprule
\textbf{Claim} & \textbf{Status} & \textbf{Reference} \\
\midrule
\endhead
$\binom{4}{3} = 4$ uniquely & Exact arithmetic & Theorem~\ref{thm:fixedpoint} \\
$K_4$ unique fixed point of $\Sop$ & Analytical theorem & Theorem~\ref{thm:fixedpoint} \\
$|\Sop^k(K_n)|$ strictly increasing for $n \geq 5$ & Analytical theorem & Theorem~\ref{thm:growth} \\
$\Rop^1$ and $\Rop^2$ structurally distinct for $n \geq 5$ & Analytical theorem & Theorem~\ref{thm:separation} \\
$\nlife < \ncons$ & Analytical theorem & Theorem~\ref{thm:separation} \\
$\nmax < \infty$ & Analytical theorem & Theorem~\ref{thm:nmax} \\
$\nlife < \ncons < \nmax$ & Analytical theorem & Theorem~\ref{thm:hierarchy} \\
$\deltastar = 1/2$ & Comp.\ result & DL01, Script~3 \\
\midrule
$\nlife$ is finite (existence) & Open problem & DL03 \\
Numerical value of $\nlife$ & Open problem & DL03 \\
Numerical value of $\ncons$ & Open problem & DL04--05 \\
Numerical value of $\nmax$ & Open problem & Computable \\
Gain function of $\Rop^1_{\mathrm{active}}$ & Open problem & DL03 \\
\midrule
Life as $\Rop^1_{\mathrm{active}}$ at $\nlife$ & Conj.\ (struct.) & PDL-V Part~1 \\
Consciousness as $\Rop^2_{\mathrm{active}}$ at $\ncons$ & Interp.\ remark & DL01, §7 \\
\bottomrule
\end{longtable}
\end{center}

% ============================================================
\section{Conclusion}
\label{sec:conclusion}
% ============================================================

Three structural results have been established in this document, all as analytical
consequences of C1--C4 and the single parameter $\kap = 1/4$.

The first is the uniqueness of $\Kfour$ as a fixed point of $\Sop$. This is not a
coincidence of the PDL framework --- it is the unique positive integer solution of
$\binom{n}{3} = n$. It explains why the self-representation hierarchy is rooted at
the quantum level, and why $n = 4$ is the only level where life and consciousness
cannot yet be distinguished.

The second is the separation of the life and consciousness thresholds. For any $n
\geq 5$, the conditions $\Rop^1$ and $\Rop^2$ are structurally different: the target
of the required homomorphism grows exponentially between the two levels. There
necessarily exist structures that are alive without being conscious, and the gap
between the two conditions grows with the complexity of the structure.

The third is the finiteness of the thermodynamic ceiling. The rarity law
$f(n) = 2^{n+1-\binom{n}{2}}$, derived from C1--C4 in DL01, decreases doubly
exponentially. There exists a finite level beyond which no coherent closure can be
thermodynamically maintained.

Together, these three results establish the following as a theorem of C1--C4: the
hierarchy $n = 4 < \nlife < \ncons < \nmax < \infty$ is logically inevitable. Life
and consciousness are not mysteries added to physics from outside --- they are required
levels in a hierarchy that the four axioms force. The present document proves that
these levels exist, are ordered, and are bounded. The programme of DL03--DL05 is to
determine where they are.

\vspace{1em}
\begin{center}
\textit{The axioms do not choose where life begins. They only ensure that it must.}
\end{center}

\clearpage

% ============================================================
\bibliographystyle{unsrt}
\bibliography{DL02_references}
% ============================================================

\end{document}
